\section*{第 \ref{ch:g11:deriv} 章 --- 导数：入门教程}

\begin{solution}{exo:g11:deriv:1}
对 $f(x) = x^2$，取 $a = 2$：
\[
\frac{(2+h)^2 - 4}{h} = \frac{4h + h^2}{h} = 4 + h
\xrightarrow[h\to0]{} 4,
\]
于是 $f'(2) = 4$。对 $g(x) = x^2 + 3x$，取 $a = 1$，$g(1) = 4$，且
\[
\frac{(1+h)^2 + 3(1+h) - 4}{h} = \frac{5h + h^2}{h} = 5 + h
\xrightarrow[h\to0]{} 5,
\]
于是 $g'(1) = 5$。
\end{solution}

\begin{solution}{exo:g11:deriv:2}
$f'(x) = 12x^2 - 4x + 1$.

$g'(x) = \frac{5}{2\sqrt x} - \frac{3}{x^2}$（在 $\intoo{0}{+\infty}$ 上）。

$h'(x) = 2(x^2 - 3) + (2x+1)(2x) = 2x^2 - 6 + 4x^2 + 2x
= 6x^2 + 2x - 6$（由乘积法则）。
\end{solution}

\begin{solution}{exo:g11:deriv:3}
$f(x) = x^2 - 3x$，$a = 2$：$f(2) = -2$，$f'(x) = 2x - 3$，故 $f'(2) = 1$；
切线 $y = -2 + (x - 2) = x - 4$。

$f(x) = \frac1x$，$a = 1$：$f(1) = 1$，$f'(1) = -1$；切线
$y = 1 - (x - 1) = 2 - x$。

$f(x) = \sqrt x$，$a = 4$：$f(4) = 2$，$f'(4) = \frac{1}{2\sqrt4} =
\frac14$；切线 $y = 2 + \frac14(x - 4) = \frac{x}{4} + 1$。
\end{solution}

\begin{solution}{exo:g11:deriv:4}
$f'(x) = \frac{(x-1) - (x+2)}{(x-1)^2} = \frac{-3}{(x-1)^2}$.

$g'(x) = \frac{2x(x^2+1) - x^2 \times 2x}{(x^2+1)^2}
= \frac{2x}{(x^2+1)^2}$.

$h'(x) = -\frac{2x+1}{(x^2+x+1)^2}$（由商的求导法则，取 $u = 1$）。
\end{solution}

\begin{solution}{exo:g11:deriv:5}
$f'(x) = 3x^2 - 12x + 9 = 3(x^2 - 4x + 3) = 3(x-1)(x-3)$：
在 $\intcc{1}{3}$ 之外为正，在其内为负。故 $f$ 在
$\intoc{-\infty}{1}$ 上递增，在 $\intcc{1}{3}$ 上递减，在
$\intco{3}{+\infty}$ 上递增，有极大值 $f(1) = 1 - 6 + 9 - 2 = 2$，
极小值为 $f(3) = 27 - 54 + 27 - 2 = -2$。
\end{solution}

\begin{solution}{exo:g11:deriv:6}
在 $\intoo{-1}{+\infty}$ 上：
\[
f'(x) = \frac{2x(x+1) - (x^2+3)}{(x+1)^2}
= \frac{x^2 + 2x - 3}{(x+1)^2}
= \frac{(x+3)(x-1)}{(x+1)^2}.
\]
当 $x > -1$ 时，$x + 3$ 与 $(x+1)^2$ 均为正，故 $f'(x)$ 与
$x - 1$ 同号：$f$ 在 $\intoc{-1}{1}$ 上递减，在 $\intco{1}{+\infty}$ 上递增，
且最小值为 $f(1) = \frac{4}{2} = 2$。
\end{solution}

\begin{solution}{exo:g11:deriv:7}
$f'(x) = 2x - 4$.

\emph{1.} 水平切线：令 $f'(x) = 0$，得 $x = 2$；切点为
$(2, f(2)) = (2, 1)$（抛物线的顶点）。

\emph{2.} 与 $y = 2x + 1$ 平行，即斜率为 $2$：$2x - 4 = 2$，故
$x = 3$，切点为 $(3, f(3)) = (3, 2)$。
\end{solution}

\begin{solution}{exo:g11:deriv:8}
设宽为 $x$（有围栏的两边），长为 $60 - 2x$：
$A(x) = x(60 - 2x) = 60x - 2x^2$，定义在 $\intoo{0}{30}$ 上。于是
$A'(x) = 60 - 4x$，在 $x = 15$ 之前为正，之后为负，故面积
在 $x = 15$ 时最大，此时围成一块 $15 \times 30$ m 的围场，面积为
$450$ m$^2$。由\cref{ch:g11:quad}中的顶点公式
$\alpha = -\frac{60}{2 \times (-2)} = 15$ 可得同样结果，无需微积分。
\end{solution}

\begin{solution}{exo:g11:deriv:9}
设 $f(x) = \frac{x+1}{2} - \sqrt x$，定义在 $\intoo{0}{+\infty}$ 上。则
\[
f'(x) = \frac12 - \frac{1}{2\sqrt x} = \frac{\sqrt x - 1}{2\sqrt x},
\]
其符号与 $\sqrt x - 1$ 相同：在 $\intoo{0}{1}$ 上为负，
在 $\intoo{1}{+\infty}$ 上为正。故 $f$ 先递减后递增，最小值为
$f(1) = 1 - 1 = 0$。因此 $f(x) \geq 0$ 对所有 $x > 0$ 成立，即
$\sqrt x \leq \frac{x+1}{2}$，且仅当 $x = 1$ 时取等号。
\end{solution}

\begin{solution}{exo:g11:deriv:10}
由 $\pi r^2 h = 500$，得 $h = \frac{500}{\pi r^2}$，故
\[
S(r) = 2\pi r^2 + 2\pi r \cdot \frac{500}{\pi r^2}
= 2\pi r^2 + \frac{1000}{r},
\qquad
S'(r) = 4\pi r - \frac{1000}{r^2}.
\]
当 $S'(r) = 0$ 时，$4\pi r^3 = 1000$，即 $r^3 = \frac{250}{\pi}$，故
$r = \sqrt[3]{250/\pi}$（$\approx 4.3$ cm）。由于 $S'(r) =
\frac{4\pi r^3 - 1000}{r^2}$ 在此值之前为负，之后为正，
故此点为极小值点。
\end{solution}

\begin{solution}{exo:g11:deriv:11}
曲线 $y = x^2$ 在横坐标为 $a$ 的点处的切线为
$y = a^2 + 2a(x - a) = 2ax - a^2$。该切线经过 $P(1, -3)$ 的条件是
$-3 = 2a - a^2$，即 $a^2 - 2a - 3 = 0$，其根为 $a = 3$ 与
$a = -1$（$\Delta = 16$）。因此过
$P$ 恰有 \emph{两}条切线：
\[
y = 6x - 9 \quad (a = 3),
\qquad
y = -2x - 1 \quad (a = -1).
\]
\end{solution}

\begin{solution}{exo:g11:deriv:12}
由\cref{ex:g11:deriv:study}可知，$f$ 先递增至局部极大值
$f(-1) = 3$，递减至局部极小值 $f(1) = -1$，然后
又递增；当 $x$ 取大的正值（负值）时，$x^3$ 起主导作用，$f$
可取任意大（或任意小）的值。在变化表上考察水平直线
$y = k$ 与曲线的交点：
\begin{itemize}
  \item 当 $k > 3$ 或 $k < -1$ 时，直线与曲线只有一个交点
        --- 共 $1$ 个解；
  \item 当 $k = 3$ 或 $k = -1$ 时，直线与曲线在一个极值点处相切，
        并与另一支相交 --- 共 $2$ 个解；
  \item 当 $-1 < k < 3$ 时，直线与三支曲线都相交 --- 共 $3$
        个解。
\end{itemize}
\end{solution}

\begin{solution}{pb:g11:deriv:1}
\textbf{1.} $f'(x) = 6x - 5$；$g'(x) = 3x^2 - 12$。由导数定义，
$\frac{(1 + h)^2 - 1}{h} = 2 + h \to 2$，即
$x^2$ 在 $1$ 处的导数为 $2$。

\textbf{2.} $g(2) = 8 - 24 = -16$ 且 $g'(2) = 12 - 12 = 0$：
切线为水平直线 $y = -16$。水平切线
表明该点为驻点 --- 此处为局部极小值，
与第 5 题的变化表相符。

\textbf{3.} 切线斜率为 $2a$（切点为 $(a, a^2)$）：
$y = a^2 + 2a(x - a) = 2ax - a^2$.

\textbf{4.} 令 $x = 0$，得 $y = -a^2$：切线与纵轴相交于
原点下方深度为 $a^2$ 处 --- 恰与 $P$ 高出零点的距离相等，
即位于零点\emph{下方}。作图方法：要画出
高度为 $h$ 的点处的切线，只需将该点与
原点下方深度为 $h$ 的纵轴上的点相连。（无需求导，
作图时即可完成。）

\textbf{5.} $g'(x) = 3(x^2 - 4)$：
在 $\intoo{-\infty}{-2}$ 上为正，在 $\intoo{-2}{2}$ 上为负，
在其后为正。$g$ 先递增至局部极大值 $g(-2) = 16$，
再递减至局部极小值 $g(2) = -16$，然后又递增。

\textbf{6.} 点 $F\left(0, \frac14\right)$ 与 $T(0, -a^2)$：
$FT = \frac14 + a^2$.

\textbf{7.} $FP = a^2 + \frac14 = FT$：三角形在 $F$ 处为等腰，
故两底角相等：$\widehat{FTP} = \widehat{FPT}$。

\textbf{8.} 入射到 $P$ 点的光线竖直，故与
直线 $(TF)$（即 $y$ 轴）平行。切线 $(TP)$ 穿过
这两条平行线，故光线与切线的夹角
等于 $\widehat{FTP}$（内错角） --- 而它又等于
$\widehat{FPT}$，即 $[PF]$ 与切线的夹角。
因此切线在 $P$ 点两侧的角相等。

\textbf{9.} 由反射定律，竖直射到
$P$ 点的光线沿与之成等角的另一侧射出 ---
由第 8 题，沿 $[PF]$ 方向射出：经过
焦点，与 $a$ 无关。反过来，从 $F$ 发出的光线
经镜面每一点反射后都平行于对称轴。碟形天线可以聚光，
车灯可以发出平行光：证毕。

\textbf{10.} 取 $a = 1$：切线为 $y = 2x - 1$，$T(0, -1)$；
$FP = \sqrt{1 + \left(1 - \frac14\right)^2} =
\sqrt{\frac{25}{16}} = \frac54$; $FT = \frac14 + 1 = \frac54$:
等腰，正如所料。

\textbf{11.} 在 $x_n$ 处的切线为
$y = f(x_n) + f'(x_n)(x - x_n)$；它与 $y = 0$ 的交点满足
$x = x_n - \frac{f(x_n)}{f'(x_n)}$.

\textbf{12.} $x_1 = 1 - \frac{-1}{2} = \frac32$;
$x_2 = \frac32 - \frac{1/4}{3} = \frac{17}{12}$;
$x_3 = \frac{577}{408}$：这些是海伦求 $\sqrt2$ 的近似值，
即\cref{pb:g10:functions:1}中的不动点迭代过程。
恒等式：$x - \frac{x^2 - 2}{2x} = \frac{2x^2 - x^2 + 2}{2x}
= \frac{x^2 + 2}{2x} = \frac12\left(x + \frac2x\right)$ ---
海伦的算法\emph{正是}牛顿法应用于
$x^2 - 2$，比牛顿早了两千年。

\textbf{13.} $x_1 = 5 - \frac{125 - 100}{75} = \frac{14}{3}
\approx 4.6667$; $x_2 \approx 4.6667 - \frac{1.63}{65.33}
\approx 4.6417$。计算器：$\sqrt[3]{100} = 4.6416\ldots$ ---
两步即得四位正确小数。

\textbf{14.} $f'(2) = 0$：起始点处的切线
水平（第 2 题），且与横轴永不相交 --- 公式
中出现除以零的情形。注意：牛顿法的初值应远离
驻点（更一般地，应足够接近所求的
根）。

\textbf{15.} 二分法每步增加一个二进制位；牛顿法
每步大致使正确数字的位数\emph{翻倍}
（本例中依次为 $1 \to 3 \to 6$）。当每步计算代价高昂 --- 数百万
个参数、实时导航等场合，翻倍的速度胜过
一点点逼近，这正是切线广泛运行于
各类芯片之中的原因。

\textbf{16.} $d^2 = 900 - w^2$，故
$S(w) = w(900 - w^2) = 900w - w^3$，其中 $0 < w < 30$。

\textbf{17.} 令 $S'(w) = 900 - 3w^2 = 0$，解得
$w = \sqrt{300} = 10\sqrt3 \approx 17.3$ 厘米（$S' > 0$ 在其前，
在其后 $< 0$，故为极大值）。于是
$d = \sqrt{600} = 10\sqrt6 \approx 24.5$ 厘米。

\textbf{18.} $\frac{d^2}{w^2} = \frac{600}{300} = 2$，故
$d = w\sqrt2$：深度是宽度的 $\sqrt2$ 倍 --- 即
木工常用的比例（用三等分直径的作图法
只需圆规即可得到，现场无需代数运算）。

\textbf{19.} 正方形梁：$w = d = \sqrt{450} \approx 21.2$
厘米，刚度 $w^3 = 450^{3/2} \approx 9\,546$。最优梁：
$S = 10\sqrt3 \times 600 = 6\,000\sqrt3 \approx 10\,392$。增益：
同一根原木的刚度大约提高 $8.9\,\%$ --- 导数的
价值足以抵偿锯木厂的投入。

\textbf{20.} 几何上，切线的等角使
抛物面成为聚光器。数值上，沿切线下滑
即可以位数翻倍的速度求解方程。
最优化上，导数为零之处找到了原木中
最强的矩形。一物三用 --- 而十二年级的
凸性还将进一步说明，
曲线位于每一条切线的哪一侧。
\end{solution}
